% sec10_4.tex — Section 10.4: Inference on Cointegrating Vectors

In the previous section, we examined whether an I(1) system is cointegrated. In
this section, we assume that the system is known to be cointegrated and that the
cointegration rank and the associated triangular representation are known.\footnote{%
In practice, we rarely have such knowledge, and the decision to entertain a system with $h$ cointegrating
relationships is usually based on the outcome of prior tests (such as those mentioned in the previous section)
designed for determining the cointegration rank. This creates a \textbf{pretest problem}, because the distribution of the
estimated cointegrating vectors does not take into account the uncertainty about the cointegration rank. This issue
has not been studied extensively.}
Our
interest is to estimate the cointegrating vectors in the triangular representation and
to make inference about them. For the most part, we focus on the special case
where the cointegration rank is~1; the general case is briefly discussed at the end of
the section.

\subsection*{The SOLS Estimator}

For the $h = 1$ case, the triangular representation is
\begin{equation}
  \begin{cases}
    y_{1t} = \mu + \boldsymbol{\gamma}' \mathbf{y}_{2t} + z_t^* \\[3pt]
    \underset{((n-1) \times 1)}{\Delta \mathbf{y}_{2t}}
      = \boldsymbol{\delta}_2 + \mathbf{u}_{2t}
  \end{cases}, \quad
  \begin{bmatrix}
    \underset{(1 \times 1)}{z_t^*} \\[3pt]
    \underset{((n-1) \times 1)}{\mathbf{u}_{2t}}
  \end{bmatrix}
  = \underset{(n \times n)}{\boldsymbol{\Psi}^*(L)} \;
    \underset{(n \times 1)}{\boldsymbol{\varepsilon}_t},
  \label{eq:10.4.1} \tag{10.4.1}
\end{equation}
where $\mathbf{y}_{2t}$ is not cointegrated. (This is just reproducing (10.2.3) and (10.2.5) with
(10.2.39) for $h = 1$.) So the regression (10.3.1) is now the cointegrating regression.
Therefore, there exists a unique $(n - 1)$-dimensional vector $\boldsymbol{\gamma}$ such that
$y_{1t} - \hat{\boldsymbol{\gamma}}'\mathbf{y}_{2t}$
is a stationary process ($z_t^*$) plus some time-invariant random variable ($\mu$) when
$\hat{\boldsymbol{\gamma}} = \boldsymbol{\gamma}$ and has a stochastic trend when
$\hat{\boldsymbol{\gamma}} \neq \boldsymbol{\gamma}$. This suggests that the OLS estimate
$(\hat{\mu}, \hat{\boldsymbol{\gamma}})$ of the coefficient vector, which minimizes the sum of squared residuals,
is consistent. In fact, as was shown by Phillips and Durlauf (1986) and Stock
(1987) for the case where $\boldsymbol{\delta}_2$ in (10.4.1) ($= \E(\Delta\mathbf{y}_{2t})$) is zero and by Hansen (1992a,
Theorem~1(b)) for the $\boldsymbol{\delta}_2 \neq \mathbf{0}$ case, the OLS estimate $\hat{\boldsymbol{\gamma}}$ is superconsistent for the
cointegrating vector $\boldsymbol{\gamma}$, with the sampling error converging to 0 at a rate faster than
the usual rate of $\sqrt{T}$.\footnote{%
The speed of convergence is $T$ if $\boldsymbol{\delta}_2 = \mathbf{0}$, and $T^{3/2}$ if $\boldsymbol{\delta}_2 \neq \mathbf{0}$. All the studies cited here assume that $\mu$
is a fixed constant. The same conclusion should hold even if $\mu$ is treated as random. As mentioned in Section
10.1 (see the paragraph right below (10.1.21)), strictly speaking, $\mu + z_t^*$ is not stationary even if $z_t^*$ is stationary.
However, since $\mu$ and $z_t^*$ are asymptotically independent as $t \to \infty$, $\mu + z_t^*$ is asymptotically stationary, which
is all that is needed for the asymptotic results here.}
Also, the $R^2$ converges to unity.\footnote{%
For a statement of the conditions under which these results hold, see Hamilton (1994, Proposition~19.2).
Those conditions are restrictions on $\boldsymbol{\Psi}^*(L)\boldsymbol{\varepsilon}_t$ in (10.4.1).}
The OLS estimator of $\boldsymbol{\gamma}$
from the cointegrating regression will be referred to as the \textbf{``static'' OLS (SOLS)}
estimator of the cointegrating vector. Since the SOLS estimator is consistent, the
residuals converge to a zero-mean stationary process. Thus, if a univariate unit-root
test such as the ADF test is applied to the residuals, the test will reject the I(1)
null in large samples. This is why the residual-based test of the previous section is
consistent against cointegration.

This fact --- that the OLS coefficient estimates are consistent when the regres\-sors
are I(1) and not cointegrated --- is a remarkable result, in sharp contrast to the
case of stationary regressors. To appreciate the contrast, remember from Chap\-ter~3
what it takes to obtain a consistent estimator of $\boldsymbol{\gamma}$ when the regressors $\mathbf{y}_{2t}$
are stationary: for the OLS estimator to be consistent, the error term $z_t^*$ has to be
uncorrelated with $\mathbf{y}_{2t}$; otherwise we need instrumental variables for $\mathbf{y}_{2t}$. In contrast,
if $y_{1t}$ is cointegrated with $\mathbf{y}_{2t}$ and if the I(1) regressors $\mathbf{y}_{2t}$ are not cointegrated, as
here, then we do not have to worry about the simultaneity bias, at least in large
samples, even though the error term $z_t^*$ and the I(1) regressors are correlated.\footnote{%
If $\mathbf{y}_{2t}$ is cointegrated, then $h \geq 2$ and the regression (10.3.1)---with $n - 1$ regressors---is no longer a
cointegrating regression (a cointegrating regression should have $n - h$ regressors, see the triangular representation
(10.2.8)). This case can be handled by the general methodology presented by Sims, Stock, and Watson (1990).
See Hamilton (1994, Chapter~18) or Watson (1994, Section~2) for an accessible exposition of the methodology.}

In finite samples, however, the bias of the SOLS estimator (the difference
between the expected value of the estimator and the true value) can be large, as
noted by Banerjee et al.\ (1986) and Stock (1987). Another shortcoming of SOLS
is that the asymptotic distribution of the associated $t$ value depends on nuisance
parameters (which are the coefficients in $\boldsymbol{\Psi}^*(L)$ in (10.4.1)), so it is difficult to do
inference. Later in this section we will introduce another estimator (to be referred
to as the ``DOLS'' estimator) which is efficient and whose associated test statistics
(such as the $t$- and Wald statistics) have conventional asymptotic distributions.

\subsection*{The Bivariate Example}

To be clear about the source and nature of the correlation between the error term
and the I(1) regressors and also to pave the way for the introduction of the DOLS
estimator, consider the bivariate version of (10.4.1). For the time being, assume
$\boldsymbol{\Psi}^*(L) = \boldsymbol{\Psi}_0^*$ (so there is no serial correlation in $(z_t^*, u_{2t})$),
$\boldsymbol{\delta}_2 = 0$, and $\mu = 0$. The
triangular representation can be written as
\begin{equation}
  \begin{cases}
    y_{1t} = \gamma y_{2t} + z_t^* \\
    \Delta y_{2t} = u_{2t}
  \end{cases}, \quad
  \begin{bmatrix} z_t^* \\ u_{2t} \end{bmatrix}
  = \boldsymbol{\Psi}_0^* \, \boldsymbol{\varepsilon}_t.
  \label{eq:10.4.2} \tag{10.4.2}
\end{equation}
The error term $z_t^*$ and the I(1) regressor $y_{2t}$ in the cointegrating regression in
(10.4.2) can be correlated because
\begin{align}
  \text{Cov}(y_{2t}, z_t^*)
  &= \text{Cov}(y_{2,0} + \Delta y_{2,1} + \Delta y_{2,2} + \cdots + \Delta y_{2t},\; z_t^*)
  \notag \\
  &= \text{Cov}(\Delta y_{2,1} + \Delta y_{2,2} + \cdots + \Delta y_{2t},\; z_t^*)
    \quad (\text{since } \text{Cov}(y_{2,0}, z_t^*) = 0)
  \notag \\
  &= \text{Cov}(u_{2,1} + u_{2,2} + \cdots + u_{2t},\; z_t^*)
    \quad (\text{since } \Delta y_{2t} = u_{2t})
  \notag \\
  &= \text{Cov}(u_{2t}, z_t^*)
    \quad (\text{since } (u_{2t}, z_t^*) \text{ is i.i.d.}).
  \label{eq:10.4.3} \tag{10.4.3}
\end{align}
Since $\boldsymbol{\Psi}_0^*$ and $\boldsymbol{\Omega}$ are not restricted to be diagonal, $z_t^*$ and $u_{2t}$ can be contemporane\-ously
correlated.

To isolate this possible correlation, consider the least squares projection of $z_t^*$
on a constant and $u_{2t}$ ($= \Delta y_{2t}$). Recall from Section~2.9 that
$\tilde{\E}^*(y \mid 1, x) = [\E(y) - \beta_0 \E(x)] + \beta_0 x$
where $\beta_0 = \text{Cov}(x, y)/\Var(x)$ and that the least squares
projection error, $y - \tilde{\E}^*(y \mid 1, x)$, has mean zero and is uncorrelated with $x$. Here,
the mean is zero for both $z_t^*$ and $u_{2t}$. So
$\tilde{\E}^*(z_t^* \mid 1, u_{2t}) = \beta_0 u_{2t}$ and, if we denote
the least squares projection error by $v_t \equiv z_t^* - \tilde{\E}^*(z_t^* \mid 1, u_{2t})$, we have
\begin{gather}
  z_t^* = \beta_0 u_{2t} + v_t,
  \notag \\
  \E(v_t) = 0, \quad
  \text{Cov}(u_{2t}, v_t) = 0, \quad
  \text{Cov}(\Delta y_{2t}, v_t) = 0.
  \label{eq:10.4.4} \tag{10.4.4}
\end{gather}
Substituting this into the cointegrating regression, we obtain
\begin{equation}
  y_{1t} = \gamma y_{2t} + \beta_0 \Delta y_{2t} + v_t.
  \label{eq:10.4.5} \tag{10.4.5}
\end{equation}
This regression will be referred to as the \textbf{augmented cointegrating regression}.
Now we show that the I(1) regressor $y_{2t}$ is \emph{strictly exogenous} in that
$\text{Cov}(y_{2s}, v_t) = 0$ for all $t, s$. By construction,
$\text{Cov}(\Delta y_{2t}, v_t) = 0$. For $\text{Cov}(\Delta y_{2s}, v_t)$ for $t \neq s$,
\begin{align}
  \text{Cov}(\Delta y_{2s}, v_t)
  &= \text{Cov}(u_{2s}, v_t)
   = \text{Cov}(u_{2s}, z_t^* - \beta_0 u_{2t})
  \notag \\
  &= \text{Cov}(u_{2s}, z_t^*) - \beta_0 \,\text{Cov}(u_{2s}, u_{2t}) = 0,
  \label{eq:10.4.6} \tag{10.4.6}
\end{align}
because $(z_t^*, u_{2t})$ is i.i.d. So $\Delta y_{2t}$ is strictly exogenous. The strict exogeneity of
$\Delta y_{2t}$ implies the same for $y_{2t}$, because
\begin{align}
  \text{Cov}(y_{2s}, v_t)
  &= \text{Cov}(y_{2,0} + \Delta y_{2,1} + \Delta y_{2,2} + \cdots + \Delta y_{2s},\; v_t)
  \notag \\
  &= 0 \quad \text{for all } t, s \text{ by (10.4.6)}.
  \label{eq:10.4.7} \tag{10.4.7}
\end{align}

\subsection*{Continuing with the Bivariate Example}

To recapitulate, we have shown that, in the augmented cointegrating regression
(10.4.5), the I(1) regressor is strictly exogenous if $\boldsymbol{\Psi}^*(L) = \boldsymbol{\Psi}_0^*$. Let
$(\hat{\gamma}, \hat{\beta}_0)$ be
the OLS coefficient estimate of $(\gamma, \beta_0)$ from the augmented cointegrating regres\-sion.
(This $\hat{\gamma}$ should be distinguished from the SOLS estimator of $\gamma$ in the coin\-tegrating
regression in (10.4.2).) This regression, (10.4.5), is very similar to the
augmented autoregression (9.4.6) in that one of the two regressors is zero-mean
I(0) and the other is driftless I(1). We have shown for the augmented
autoregression that the ``$\mathbf{X}'\mathbf{X}$'' matrix, if properly scaled by $T$ and $\sqrt{T}$, is asymptotically
diagonal, so the existence of I(0) regressors can be ignored for the purpose of
deriving the limiting distribution of $\hat{\gamma}$, the OLS estimator of the coefficient of the
I(1) regressor. The same is true here, and the same argument exploiting the asymp\-totic
diagonality of the suitably scaled ``$\mathbf{X}'\mathbf{X}$'' matrix shows that the usual $t$-value
for the hypothesis that $\gamma = \gamma_0$ is asymptotically equivalent to
\begin{equation}
  \tilde{t}
  = \frac{\displaystyle \frac{1}{T} \sum_{t=1}^{T} y_{2t} v_t}
         {\displaystyle \sqrt{\frac{\sigma_v^2}{T^2}
           \sum_{t=1}^{T} (y_{2t})^2}},
  \label{eq:10.4.8} \tag{10.4.8}
\end{equation}
where $\sigma_v^2$ is the variance of $v_t$. That is, the difference between the usual $t$-value
and $\tilde{t}$ in (10.4.8) converges to zero in probability, so the limiting distribution of $t$
and that of $\tilde{t}$ are the same.

In the augmented autoregression used for the ADF test, the limiting distribu\-tion
of the ADF $t$-statistic was nonstandard (it is the Dickey-Fuller $t$-distribution).
In contrast, the asymptotic distribution of $\tilde{t}$ (and hence that of the usual $t$-value)
is \emph{standard normal}. To see why, first recall that the I(1) regressor $y_{2t}$ is strictly
exogenous in that $\text{Cov}(y_{2s}, v_t) = 0$ for all $t, s$. To develop intuition, \emph{temporar\-ily}
assume that $(z_t^*, u_{2t})$ is jointly normal. Then $y_{2s}$ and $v_t$ are independent for
all $(s, t)$, not just uncorrelated. Consequently, the distribution of $v_t$ conditional on
$(y_{2,1}, y_{2,2}, \ldots, y_{2T})$ is the same as its unconditional distribution, which is $N(0, \sigma_v^2)$.
So the distribution of the numerator of (10.4.8) conditional on $(y_{2,1}, y_{2,2}, \ldots,
y_{2T})$ is
\[
  N\biggl(0, \; \frac{\sigma_v^2}{T^2} \sum_{t=1}^{T} (y_{2t})^2\biggr).
\]
But the standard deviation of this normal distribution equals the denominator of
$\tilde{t}$, so
\[
  \tilde{t} \mid (y_{2,1}, y_{2,2}, \ldots, y_{2T}) \sim N(0, 1).
\]
Since this conditional distribution of $\tilde{t}$ does not depend on $(y_{2,1}, y_{2,2}, \ldots, y_{2T})$, the
\emph{un}conditional distribution of $\tilde{t}$, too, is $N(0, 1)$. (This type of argument is not new to
you; see the paragraph containing (1.4.3) of Chapter~1.) Recall from Chapter~2 that,
even if the normality assumption is dropped, the distribution of the usual $t$-value
is standard normal, albeit asymptotically, in large samples. The same conclusion
is true here: the limiting distribution of $\tilde{t}$ is $N(0, 1)$ even if $(z_t^*, u_{2t})$ is not jointly
normal. Proving this requires an argument different from the type used in Chapter~2,
because $y_{2t}$ here has a stochastic trend. We will not give a formal proof here; just
an executive summary. It consists of two parts. The first is to derive the limiting
distribution of the numerator of (10.4.8) (which can be done using a result stated
in, e.g., Proposition~18.1(e) of Hamilton, 1994, or Lemma~2.3(c) of Watson, 1994).
This distribution is nonstandard. The second is to show that normalization (10.4.8)
converts the nonstandard distribution into a normal distribution (Lemma~5.1 of Park
and Phillips, 1988).

This example of a bivariate I(1) system is special in several respects: (a) there
is no serial correlation in the error process $(z_t^*, u_{2t})$, (b) the I(1) regressor $y_{2t}$ is a
scalar, (c) $y_{2t}$ has no drift, and (d) $\mu = 0$. Of these, relaxing (a) requires some
thoughts.

\subsection*{Allowing for Serial Correlation}

If it is not the case that $\boldsymbol{\Psi}^*(L) = \boldsymbol{\Psi}_0^*$ as in (10.4.2), then $(z_t^*, u_{2t})$ is serially corre\-lated.
Consequently, the I(1) regressor $y_{2t}$ in the augmented cointegrating regres\-sion
(10.4.5) is no longer strictly exogenous because $\text{Cov}(\Delta y_{2s}, v_t)$, while still zero
for $t = s$, is no longer guaranteed to be zero for $t \neq s$. To remove this correlation,
consider the least squares projection of $z_t^*$ on the current, past, and future values of
$u_2$, not just on the current value of $u_2$. Noting that $\Delta y_{2t} = u_{2t}$, the projection can
be written as
\begin{equation}
  z_t^* = \boldsymbol{\beta}(L) \Delta \mathbf{y}_{2t} + v_t, \quad
  \boldsymbol{\beta}(L) \equiv \sum_{j=-\infty}^{\infty} \boldsymbol{\beta}_j L^j.
  \label{eq:10.4.9} \tag{10.4.9}
\end{equation}
(This $v_t$ is different from the $v_t$ in (10.4.5).) By construction, $\E(v_t) = 0$ and
$\text{Cov}(\Delta y_{2s}, v_t) = 0$ for all $t, s$. The two-sided filter $\boldsymbol{\beta}(L)$ can be of infinite order,
but suppose---for now---that it is not and $\beta_j = 0$ for $|j| > p$. Thus,
\begin{align}
  z_t^*
  &= \boldsymbol{\beta}_0 \Delta \mathbf{y}_{2t}
    + \boldsymbol{\beta}_{-1} \Delta \mathbf{y}_{2,t+1} + \cdots
    + \boldsymbol{\beta}_{-p} \Delta \mathbf{y}_{2,t+p}
  \notag \\
  &\quad + \boldsymbol{\beta}_1 \Delta \mathbf{y}_{2,t-1} + \cdots
    + \boldsymbol{\beta}_p \Delta \mathbf{y}_{2,t-p} + v_t.
  \label{eq:10.4.10} \tag{10.4.10}
\end{align}
Substituting this into the cointegrating regression in (10.4.2), we obtain the (vastly)
augmented cointegrating regression
\begin{align}
  y_{1t}
  &= \gamma y_{2t}
    + \boldsymbol{\beta}_0 \Delta \mathbf{y}_{2t}
    + \boldsymbol{\beta}_{-1} \Delta \mathbf{y}_{2,t+1} + \cdots
    + \boldsymbol{\beta}_{-p} \Delta \mathbf{y}_{2,t+p}
  \notag \\
  &\quad + \boldsymbol{\beta}_1 \Delta \mathbf{y}_{2,t-1} + \cdots
    + \boldsymbol{\beta}_p \Delta \mathbf{y}_{2,t-p} + v_t.
  \label{eq:10.4.11} \tag{10.4.11}
\end{align}
Since $\text{Cov}(\Delta y_{2s}, v_t) = 0$ for all $t$ and $s$, $\Delta y_{2t}$ is strictly exogenous, which means
(as seen above) that the level regressor $y_{2t}$ too is strictly exogenous. Thus, by
including not just the current change but also past and future changes of the I(1)
regressor in the augmented cointegrating regression, we are able to maintain the
strict exogeneity of $y_{2t}$. The OLS estimator of the cointegrating vector $\boldsymbol{\gamma}$ based
on this augmented cointegrating regression is referred to as the \textbf{``dynamic'' OLS
(DOLS)}, to distinguish it from the SOLS estimator based on the cointegrating
regression without changes in $y_2$.

There are $2 + 2p$ regressors, the first of which is driftless I(1) and the rest
zero-mean I(0). As before, with suitable scaling by $T$ and $\sqrt{T}$, the I(1) regressor
is asymptotically uncorrelated with the $2p + 1$ zero-mean I(0) regressors, so that
the suitably scaled $\mathbf{X}'\mathbf{X}$ matrix is again block diagonal and the zero-mean I(0)
regressors can be ignored for the purpose of deriving the limiting distribution of
the DOLS estimate of $\gamma$. Therefore, the expression for the random variable that is
asymptotically equivalent to the $t$-value for $\gamma = \gamma_0$ is again given by (10.4.8), the
expression derived for the case where $(z_t^*, u_{2t})$ is serially uncorrelated.

However, when $(z_t^*, u_{2t})$ is serially correlated, the derivation of the asymp\-totic
distribution of (10.4.8) is not the same as when $(z_t^*, u_{2t})$ is serially uncorre\-lated,
because the error term $v_t$ can now be serially correlated; the projection
(10.4.9), while eliminating the correlation between $\Delta y_{2s}$ and $v_t$ for all $s, t$, does
not remove its own serial correlation in $v_t$. To examine the asymptotic distribution,
let $\mathbf{V}$ ($T \times T$) be the autocovariance matrix of $T$ successive values of $v_t$ and
$\lambda_v^2$ be
the long-run variance of $v_t$. Also, let $\mathbf{X}$ in the rest of this subsection be the $T \times 1$
matrix $(y_{2,1}, y_{2,2}, \ldots, y_{2T})'$. Again, to develop intuition, temporarily assume that
$(z_t^*, u_{2t})$ are jointly normal. Then, since $y_{2t}$ is strictly exogenous, the distribution
of the numerator of (10.4.8) conditional on $\mathbf{X}$ is normal with mean zero and the
conditional variance
\begin{equation}
  \frac{1}{T^2} \mathbf{X}'\mathbf{V}\mathbf{X}.
  \label{eq:10.4.12} \tag{10.4.12}
\end{equation}
The square root of this would have to replace the denominator of (10.4.8) to make
the ratio standard normal in finite samples. Fortunately, it turns out (see Corollary
2.7 of Phillips, 1988) that, in large samples, the distribution of the numerator is the
same as the distribution you would get if $v_t$ were serially uncorrelated but with the
variance of $\lambda_v^2$ rather than $\sigma_v^2$.

Therefore, all that is required to modify the ratio (10.4.8) to accommodate
serial correlation in $v_t$ is to replace $\sigma_v^2$ ($\equiv \Var(v_t)$) by $\lambda_v^2$ (the long-run variance of
$v_t$), namely, if $\tilde{t}$ is given by (10.4.8) with $\sigma_v^2$ still in the denominator, its rescaled
value
\begin{equation}
  \left(\frac{\sigma_v}{\lambda_v}\right) \cdot \tilde{t}
  \label{eq:10.4.13} \tag{10.4.13}
\end{equation}
is asymptotically $N(0, 1)$! Since the OLS $t$-value for $\gamma$ is asymptotically equivalent
to $\tilde{t}$, it follows that
\begin{equation}
  \left(\frac{s}{\hat{\lambda}_v}\right) \cdot t \dto N(0, 1),
  \label{eq:10.4.14} \tag{10.4.14}
\end{equation}
where $\hat{\lambda}_v$ is some consistent estimator of $\lambda_v$ and $s$ is the usual OLS standard error
of regression (it is easy to show that $s$ is consistent for $\sigma_v$). Put differently, if we
rescale the usual standard error of $\hat{\gamma}$ (the OLS estimate of the $y_{2t}$ coefficient in
(10.4.11)) as
\begin{equation}
  \text{rescaled standard error}
  = \left(\frac{\hat{\lambda}_v}{s}\right) \times \text{usual standard error},
  \label{eq:10.4.15} \tag{10.4.15}
\end{equation}
then the $t$-value based on this rescaled standard error is asymptotically $N(0, 1)$.

The foregoing argument is applicable only to the coefficient of the I(1) regres\-sor,
so the $t$-values for $\boldsymbol{\beta}$'s in (10.4.11), even when their standard errors are rescaled
as just described, are not necessarily $N(0, 1)$ in large samples.

Since the regressors are strictly exogenous, the discussion in Sections~1.6 and
6.7 suggests that the GLS might be applicable. However, the fact noted above that
$\mathbf{X}'\mathbf{V}\mathbf{X}$ behaves asymptotically like $\lambda_v\mathbf{X}'\mathbf{X}$ implies that the GLS estimator of
$\gamma$ (to
be referred to as the \textbf{DGLS} estimator) is asymptotically equivalent to the DOLS
estimator (or more precisely, $T$ times the difference converges to 0 in probability)
(see Phillips, 1988, and Phillips and Park, 1988). Therefore, there is no efficiency
gain from correcting for the serial correlation in the error term $v_t$ by GLS.

A consistent estimate of $\lambda_v$ is easy to obtain. Recall that in Section~6.6 we
used the VARHAC procedure to estimate the long-run variance matrix of a vector
process. The same procedure can be applied to the present case of a scalar error
process. Consider fitting an AR($p$) process to the residuals, $\hat{v}_t$, from the augmented
cointegrating regression
\begin{equation}
  \hat{v}_t = \phi_1 \hat{v}_{t-1} + \phi_2 \hat{v}_{t-2}
    + \cdots + \phi_p \hat{v}_{t-p} + e_t
  \quad (t = p + 1, \ldots, T).
  \label{eq:10.4.16} \tag{10.4.16}
\end{equation}
(This $p$ should not be confused with the $p$ in the augmented cointegrating regres\-sion.)
Use the BIC to pick the lag length by searching over the possible lag lengths
of $0, 1, \ldots, T^{1/3}$. Using the relation (6.2.21) on page~385 and noting that the long-run
variance is the value at $z = 1$ of the autocovariance-generating function, the
long-run variance $\lambda_v^2$ of $v_t$ can be estimated by
\begin{equation}
  \hat{\lambda}_v^2
  = \frac{\hat{\sigma}_e^2}{(1 - \hat{\phi}_1 - \cdots - \hat{\phi}_p)^2}, \quad
  \hat{\sigma}_e^2
  = \frac{1}{T - p} \sum_{t=p+1}^{T} \hat{e}_t^2,
  \label{eq:10.4.17} \tag{10.4.17}
\end{equation}
where $\hat{\phi}_j$ ($j = 1, 2, \ldots, p$) are the estimated AR($p$) coefficients and $\hat{e}_t$ is the
residual from the AR($p$) equation (10.4.16).

\subsection*{General Case}

We now turn to the general case of an $n$-dimensional cointegrated system with
possibly nonzero drift. We focus on the $h = 1$ case because extending it to the case
where $h > 1$ is straightforward. So the triangular representation is (10.4.1) and the
augmented cointegrating regression is
\begin{align}
  y_{1t}
  &= \mu + \boldsymbol{\gamma}' \mathbf{y}_{2t}
    + \boldsymbol{\beta}_0' \Delta \mathbf{y}_{2t}
    + \boldsymbol{\beta}_{-1}' \Delta \mathbf{y}_{2,t+1} + \cdots
    + \boldsymbol{\beta}_{-p}' \Delta \mathbf{y}_{2,t+p}
  \notag \\
  &\quad + \boldsymbol{\beta}_1' \Delta \mathbf{y}_{2,t-1} + \cdots
    + \boldsymbol{\beta}_p' \Delta \mathbf{y}_{2,t-p} + v_t.
  \label{eq:10.4.18} \tag{10.4.18}
\end{align}
Here, $\boldsymbol{\beta}_j$ ($j = 0, 1, 2, \ldots, p, -1, -2, \ldots, -p$) are the least squares projection
coefficients in the projection of $z_t^*$ (the error in the cointegrating regression) on the
current, past, and future values of $\Delta\mathbf{y}_2$. The DOLS estimator of the cointegrating
vector $\boldsymbol{\gamma}$ is simply the OLS estimator of $\boldsymbol{\gamma}$ from this augmented cointegrating
regression.

The following results are proved in Saikkonen (1991) and Stock and Watson
(1993).\footnote{%
These authors assume that $\mu$ is a fixed constant. However, the same conclusions would hold even if $\mu$ were
a time-invariant random variable.}

\begin{enumerate}
\item The bivariate results derived above carry over. That is, the DOLS estimator
  of $\boldsymbol{\gamma}$ is superconsistent and the properly rescaled $t$- and Wald statistics for
  hypotheses about $\boldsymbol{\gamma}$ have the conventional asymptotic distributions (standard
  normal and chi squared). The proper rescaling is to multiply the usual $t$-value
  by $(s/\hat{\lambda}_v)$ and the Wald statistics by the square of $(s/\hat{\lambda}_v)$. This simple rescaling,
  however, does not work for the $t$-value and the Wald statistic for hypotheses in\-volving
  $\mu$ or $\boldsymbol{\beta}_j$.

\item If the two-sided filter $\boldsymbol{\beta}(L)$ in the projection of $z_t^*$ on the leads and lags of
  $\Delta\mathbf{y}_{2t}$
  is infinite order, then the error term $v_t$ includes the truncation remainder,
  \[
    \sum_{j=p+1}^{\infty} \boldsymbol{\beta}_{-j}' \Delta\mathbf{y}_{2,t+j}
    + \sum_{j=p+1}^{\infty} \boldsymbol{\beta}_j' \Delta\mathbf{y}_{2,t-j}.
  \]
  All the results just mentioned above carry over, provided that $p$ in (10.4.18)
  is made to increase with the sample size $T$ at a rate slower than $T^{1/3}$. See
  Saikkonen (1991, Section~4) for a statement of required regularity conditions.

\item The estimator is efficient in some precise sense.\footnote{%
  The $t$-value is asymptotically $N(0, 1)$, but the DOLS estimator itself is not asymptotically normal. So the
  usual criterion of comparing the variances of the asymptotic distributions among asymptotically normal estima\-tors
  is not applicable here. See Saikkonen (1991, Section~3) for an appropriate definition of efficiency.}
  The estimator is asymp\-totically
  equivalent to other efficient estimators such as Johansen's maximum
  likelihood procedure based on the VECM, the Fully Modified estimator of
  Phillips and Hansen (1990), and a nonlinear least squares estimator of Phillips
  and Loretan (1991) (and also to the DGLS, as mentioned above). For all these
  efficient estimators, the $t$- and Wald statistics for hypotheses about $\boldsymbol{\gamma}$, cor\-rectly
  rescaled if necessary, have standard asymptotic distributions (see
  Watson, 1994, Section~3.4.3, for more details).
\end{enumerate}

\subsection*{Other Estimators and Finite-Sample Properties}

Stock and Watson (1993) examined the finite-sample performance of these estima\-tors
just mentioned (excluding the Phillips-Loretan estimator). For the DGPs and
the sample sizes ($T = 100$ and 300) they examined, the following results emerged.
(1)~The bias is smallest for the Johansen estimator and largest for SOLS. (2)~How\-ever,
the variance of the Johansen estimator is much larger than those for the other
efficient estimators, and DOLS has the smallest root mean square error (RMSE).
(3)~For all estimators examined, the distribution of the $t$-values (correctly rescaled
if necessary) tends to be spread out relative to $N(0, 1)$, suggesting that the null will
be rejected too often. These conclusions are broadly consistent with the assessment
of Monte Carlo studies found in Maddala and Kim (1998, Section~5.7).

\subsection*{Question for Review}

\begin{enumerate}
\item As we have emphasized, there is a similarity between augmented autoregres\-sion
  (9.4.6) on page~587 and augmented cointegrating regression (10.4.5). Then
  why is the $t$-value on the I(1) regressor in the augmented cointegrating regres\-sion
  (10.4.5) asymptotically $N(0, 1)$ while that in (9.4.6) is not?
\end{enumerate}
